% ------------------------------------
% 参考文献
%
% ⚠️ 待补：初稿全文没有一处 \cite，仓库内也没有本队整理的文献清单。
% 规范（paper-spec §三.2）禁止虚构文献，因此这里**只保留骨架**，
% 未凭猜测添加任何条目。需要编程手/论文手提供真实文献清单后逐条填入。
%
% 下列唯一一条是赛题本身，属可核实的事实性出处，不是虚构引用。
% ------------------------------------

% AI 工具使用声明：文本由编程手提供，逐字采用；仅把中文行文中的半角分号
% 改为全角（排版惯例，不改内容）。
\section{AI 工具使用声明}

本参赛队在竞赛过程中使用了 AI 工具，主要用于论文语言润色、代码调试、\LaTeX{} 排版、图形优化及建模结果整理，详细使用情况见支撑材料。问题分析、核心建模、算法设计、参数设定、结果解释和最终结论均由参赛队主导完成；对 AI 参与完成的内容，参赛队进行了逐项人工审查、程序复算和结果核验。

\clearpage
\begin{thebibliography}{99}

\bibitem{contest2026c} 全国大学生数学建模竞赛组委会. 2026 年高教社杯全国大学生数学建模竞赛 C 题：微电网电力调控优化[Z]. 2026.

% \bibitem{...} 待补：储能调度 / 随机规划（SAA）/ 电价预测方向的真实文献，
%         请提供后可核实的完整条目（作者、题名、出处、年份、页码）。

\end{thebibliography}
